% Generated from locale/tr/content/sets-functions-relations/inductive-defs-proofs/introduction.tex; do not hand-edit.


\olfileid[tr]{sfr}{ind}{int}

\olsection{Giriş}

\begin{explain}
Tümevarım, matematikte ve mantıkta yaygın biçimde kullanılan bir ispat
tekniğidir. Doğal sayılar üzerinde tümevarım yapılan matematik derslerinden
ona aşina olabilirsiniz. Matematikte ve mantıkta tümevarıma elverişli pek çok
başka nesne kümesi de vardır.

Genel olarak tümevarım evrensel bir iddiayı ispatlamayı, yani belirli türdeki
her nesnenin bir özelliği taşıdığını göstermeyi sağlar. Özellikle ``evrensel
giriş'' ispat yöntemi fazla karmaşık olduğunda yararlıdır. Tümevarım yalnızca
özel biçimlerde kurulmuş matematiksel nesneler üzerinde işler: Kümedeki her
eleman ya temel bir elemansa ya da temel elemanlardan kuruluyorsa küme
üzerinde tümevarım kullanabiliriz. Daha kesin olarak:
\end{explain}

\begin{defn}
Bir $S$ kümesi, $x \in S$ olduğunda $f(x) \in S$ de oluyorsa $f$
fonksiyonu altında \emph{kapalı}dır.
\end{defn}

\begin{explain}
Örneğin doğal sayılar ($\mathbb{N}$), $f(x) = x+1$ ardıl fonksiyonu altında
kapalıdır. $x$ bir doğal sayıysa $x+1$ de bir doğal sayıdır. Buna karşılık
doğal sayılar $f(x) = x-1$ öncül fonksiyonu altında kapalı değildir; çünkü
$0$ doğal sayı olduğu halde $-1$ değildir.
\end{explain}

\begin{defn}
Bir $P$ özelliği, $f$'nin tanım kümesindeki her $a$ nesnesi için $P(a)$,
$P(f(a))$ sonucunu gerektiriyorsa $f(x)$ fonksiyonu altında \emph{korunur}
(yani $a$, $P$ özelliğini taşıyorsa $f(a)$ da taşır).
\end{defn}

\begin{explain}
``Çift sayı olma'' özelliği $f(x) = x+2$ fonksiyonu altında korunur; çünkü
$x$ çiftse $x+2$ de çifttir. Çift sayı olma özelliği ardıl fonksiyonu altında
korunmaz; çünkü $x$ çiftse $x+1$ tektir.

Tümevarım doğal sayılar üzerinde işler; çünkü her doğal sayıya, ardıl
fonksiyonunu sonlu sayıda uygulayarak $0$'dan ulaşılabilir. Bu, üzerinde
tümevarım yapabileceğimiz her kümenin ayırt edici özelliğidir.
\end{explain}

\begin{defn}[$\Nat$ üzerinde tümevarım]
$0$, bir $P$ özelliğini taşıyor ve $P$ ardıl fonksiyonu altında korunuyorsa
her doğal sayı $P$ özelliğini taşır.
\end{defn}

\begin{ex} İlk $n$ doğal sayının toplamı $\frac{n(n+1)}{2}$'dir. \\

\textbf{Temel durum.} İlk $0$ doğal sayının toplamı
$0=\frac{0\cdot 1}{2}$'dir.\\

\textbf{Tümevarım adımı.} Özelliğin $k$ için geçerli olduğunu varsayın.
$k+1$ için de geçerli olduğunu göstereceğiz. Başka bir deyişle $P(k)$'nin,
yani
\[ 0 + 1 + \ldots + k = \frac{k(k+1)}{2} \]
eşitliğinin geçerli olduğunu varsayın. $P(k)$ varsayımını kullanarak
$P(k+1)$'in, yani
\[ 0 + 1 + \ldots + k + (k+1) = \frac{(k+1)(k+2)}{2} \]
eşitliğinin geçerli olduğunu göstereceğiz.

\begin{align*}
0 + 1 + \ldots + k &= \frac{k(k+1)}{2} \tag{Varsayım} \\
0 + 1 + \ldots + k + (k+1) &= \frac{k(k+1)}{2} + (k+1)
\tag{Her iki tarafa $k+1$ ekleyin} \\
&= \frac{k(k+1) + 2(k+1)}{2} \\
&= \frac{k^2 + 3k +2}{2} \\
&=\frac{(k+1)(k+2)}{2}
\end{align*}
Dolayısıyla tümevarım adımı geçerlidir ve iddiayı ispatladık.
\end{ex}

\begin{explain}
!!{formula}s için temel elemanlar atomik !!{formula}s{}dir. Daha karmaşık
!!{formula}s, daha az karmaşık !!{formula}s{}in işleçlerle birleştirilmesiyle
elde edilir. Her işleç, iki !!{formula}{}ü söz konusu işleçle birleştiren yeni
bir !!{formula} döndüren bir fonksiyona karşılık gelir. Örneğin iki
!!{formula} olan $!A$ ve $!B$'yi bir tümel evetlemede birleştiren fonksiyon
$\mathcal E _\land (!A, !B) = !A \land !B$ biçiminde yazılabilir. Bunlara
\emph{formül-kurma fonksiyonları} diyebiliriz. !!{formula}s kümesi bu
fonksiyonlar altında kapalıdır: $!A$ ve $!B$ !!{formula}s olduğunda
$\lnot !A$, $!A \land !B$ vb. de !!{formula}s{}dir.
\end{explain}

\begin{defn}[!!{formula}s üzerinde tümevarım]
Her atomik !!{formula}, $P$ özelliğini taşıyor ve $P$ formül-kurma
fonksiyonları altında korunuyorsa her !!{formula}, $P$ özelliğini taşır.
\end{defn}

\begin{explain}
Bu tanım, !!{formula}s üzerinde tümevarımsal ispatlar için bir ``tarif''
önerir. Her !!{formula}{}ün $P$ özelliğini taşıdığını göstermemiz isteniyorsa
aşağıdaki adımları izleriz:\\

\textbf{Temel durum.} $!A$ atomik bir !!{formula} olsun. [\ldots]
Dolayısıyla $!A$, $P$ özelliğini taşır.

\textbf{Tümevarım adımı.} $!A$ ve $!B$, her ikisi de $P$ özelliğini taşıyan
!!{formula}s olsun.

$\lnot$ durumu: [\ldots] Dolayısıyla $\lnot !A$, $P$ özelliğini taşır.

$\land$ durumu: [\ldots] Dolayısıyla $!A \land !B$, $P$ özelliğini taşır.

$\lor$ durumu: [\ldots] Dolayısıyla $!A \lor !B$, $P$ özelliğini taşır.

$\lif$ durumu: [\ldots] Dolayısıyla $!A \lif !B$, $P$ özelliğini taşır.

$\liff$ durumu: [\ldots] Dolayısıyla $!A \liff !B$, $P$ özelliğini taşır.

$\lforall[]$ durumu: [\ldots] Dolayısıyla $\lforall[x] !A$, $P$ özelliğini taşır.

$\lexists[]$ durumu: [\ldots] Dolayısıyla $\lexists[x] !A$, $P$ özelliğini taşır. \\

Dolayısıyla her !!{formula}, $P$ özelliğini taşır.
\end{explain}

